\section{Scope, conventions, and source status}
\label{scope:sec}

\subsection{What this continuation establishes}

The task is to extend the exact-identity programme in the ProveIt
polylogarithm manuscript, not to reproduce results already supplied
by its incoming research reports. We inspected the canonical
manuscript and the six archives present in \src{docs/incoming}
at commit \src{dcd95baeb7c0e914b138bddef6829c5b1d52b726}
\cite{RepoMain,RepoEndpoint,RepoDougall,RepoTransport,RepoExact,
RepoCoincident,RepoTwisted}.
The mathematical work concentrates on four gaps explicitly left
by those sources. The following table records the resulting status.

\begin{longtable}{@{}>{\raggedright\arraybackslash}p{0.26\textwidth}>{\raggedright\arraybackslash}p{0.64\textwidth}@{}}
\toprule
Source target & Result and its precise scope\\
\midrule
\endhead
Higher spectral derivatives of the entire harmonic difference
\cite{RepoExact}
& Theorem~\ref{hn:thm:newton} and Corollary~\ref{hn:cor:jets}
give convergent formulas at every order. Theorem~\ref{hn:thm:AKbridge}
identifies all shift jets with classical multiple Arakawa--Kaneko
functions. A finite reduction of every general spectral derivative
to ordinary zeta derivatives is not asserted.\\[6pt]
Mixed Dougall jets at $b=c=1$, Question~3 of \cite{RepoTransport}
& Proposition~\ref{dg:polynomial} proves termwise collapse on the
integer residue grid. Theorems~\ref{dg:generator} and
\ref{dg:spectral} give every mixed parameter row and every
mixed spectral--Stieltjes derivative.\\[6pt]
Genuinely unequal twists, the first research question in \cite{RepoTwisted}
& Theorem~\ref{tw:tailtheorem} gives a convergent normal form at
arbitrary orders. The all-index formula \eqref{tw:alljets}
and finite scalar reductions retain the phase and endpoint
normalizations. The scalar classification is not an arithmetic
independence theorem for Fourier samples.\\[6pt]
Mixed triple collisions, Question~1 of \cite{RepoExact}
& Theorems~\ref{cl:collision} and \ref{cl:allfactor} give an
explicit triple collision expansion and the local completion for
any finite number of factors. Unqualified equality of the geometric
collision constant and the merged finite part is false. Theorem~\ref{cl:total-derivative} also resolves the one-derivative cancellation target.\\
\bottomrule
\end{longtable}

The results are proved below as analytic identities with specified
domains and normalizations. Exact finite calculations are supporting
certificates for selected algebraic steps. Numerical comparisons
test implementations and signs; they do not supply the infinite
identities or establish arithmetic independence. The supplied proofs
have not been formalized in a proof assistant.

\subsection{Classical antecedents and attribution}

The basic functions and summations are established mathematics.
We use the Hurwitz-zeta shift and derivative formulas, the
polygamma series, the Lerch transcendent, Euler's hypergeometric
transformation, and the Rogers--Dougall ${}_5F_4(1)$ summation
\cite{DLMFHurwitz,DLMFPolygamma,DLMFLerch,DLMFHypergeom,DLMFDougall}.
The multiple Arakawa--Kaneko functions were already defined in
\cite{AK1999}; related multiple and Hurwitz extensions are treated
in \cite{KT2018,CC2010}. The Hurwitz--Tornheim relation used in the
collision section has the antecedent \cite{EspinosaMoll}.
The centered resonance theorem and the entire harmonic subtraction
are results of the supplied reports. Where needed, their analytic
mechanisms are reproduced to make the new deductions reviewable
without reconstructing an entire earlier archive.

The contribution claimed here is the collection of deductions,
normal forms, convergence arguments, and corrections relative to
that inspected source set. We make no exhaustive priority claim
over the full special-functions literature. In particular, a
recognition of a classical Arakawa--Kaneko function is recorded as
an identification, not as the invention of a new special function.

\subsection{Notation and the meaning of finite part}

Write $(a)_n=\Gamma(a+n)/\Gamma(a)$ for a rising factorial,
$(a)^{\underline n}=a(a-1)\cdots(a-n+1)$ for a falling factorial,
and $[u^r]$ for coefficient extraction. The symbols
$\left\{\begin{smallmatrix}m\\j\end{smallmatrix}\right\}$ and
$\left[\begin{smallmatrix}m\\j\end{smallmatrix}\right]$ denote
Stirling numbers of the second kind and unsigned Stirling numbers
of the first kind, respectively. Empty products are $1$.
The ordinary generalized harmonic numbers are
\[
 H_n^{(r)}=\sum_{k=1}^n k^{-r},\qquad H_0^{(r)}=0,
 \qquad H_n=H_n^{(1)}.
\]
General shifts are defined explicitly in the relevant sections.
The letters $a,b,u$ are local parameters and are reintroduced
when their role changes.

We set $\psi=\Gamma'/\Gamma$, and $\psi^{(r)}$ denotes an
argument derivative. For zeta functions, $\zeta^{(k)}(s,x)$
denotes a derivative in the first, spectral variable; derivatives
in $x$ are written $\partial_x$. The generalized Stieltjes
constants use the convention
\begin{equation}\label{scope:stieltjes}
 \zeta(1+z,x)=\frac1z+
       \sum_{m=0}^{\infty}\frac{(-1)^m}{m!}\gamma_m(x)z^m,
 \qquad \gamma_m=\gamma_m(1),\quad \gamma=\gamma_0.
\end{equation}
In particular $\gamma_0(x)=-\psi(x)$.
Principal logarithms are used for powers in the right half-plane.
The phases for two-sided Fourier powers are fixed separately
in \eqref{tw:log}. A square at a complex shift is an algebraic
square, not an absolute-value square.

A series is \emph{normally convergent} on a parameter domain when
it converges absolutely and uniformly on each compact subset.
This is the hypothesis used for holomorphic differentiation of
infinite sums. The distribution-valued version means normal
convergence after pairing with each smooth test function, together
with the displayed uniform growth estimates on Fourier coefficients.

For a cutoff expansion in a positive ambient distance $\eta$,
$\FP$ retains the coefficient of $\eta^0(\log\eta)^0$ after
the divergent local terms are removed. At each singular grid
point in Section~\ref{cl:sec}, the cutoff is on the right in
the ambient coordinate $x$. This specifies a value; a change of
coordinate can change it. The symbol $\CT_\epsilon$ has the
analogous meaning for a geometric collision parameter $\epsilon$.
For multivariate meromorphic functions with poles on sums of
variables, an unqualified rectangular Laurent constant is not a
definition. A chamber, a spectral restriction, or the holomorphic
completion must first be specified.

\subsection{Reading the four proof mechanisms}

The harmonic argument first proves a compact-uniform estimate for
finite differences, then identifies their Mellin transform. This
order prevents the invalid differentiation of a formula known only
at discrete spectral values. The Dougall argument keeps the Gamma
ratios and the finite small-index terms intact until after the
subtractions have been differentiated. The unequal-twist argument
isolates finitely many modes before applying a binomial expansion,
so that the logarithm branches agree on the remaining tail.
The collision argument works with the complete local numerator of
each resonant subset, retaining finite terms as well as residues.
These are the analytic steps on which the extensions depend.
